\chapter{Matrices}

\noindent\textit{Une matrice est un tableau de nombres rangés en lignes et en colonnes.
Cet outil, au cœur de l'algèbre linéaire, permet de calculer rapidement, de résoudre des
systèmes linéaires, de décrire les transformations du plan et de modéliser des situations
(suites couplées, réseaux). C'est un chapitre spécifique du programme camerounais de la
série~C, à la jonction du calcul et de la géométrie.}
\vspace{8pt}

% ═══════════════════════════════════════════════════════════════
\section{Matrices : définitions et types}

\begin{definition}[Matrice]
Une \textbf{matrice} de \textbf{format} (ou \textbf{type}) $n\times p$ est un tableau de
$n\times p$ nombres réels rangés en $n$ \textbf{lignes} et $p$ \textbf{colonnes} :
\[
A=(a_{ij})_{\substack{1\le i\le n\\ 1\le j\le p}}=
\begin{pmatrix} a_{11} & a_{12} & \cdots & a_{1p}\\ a_{21} & a_{22} & \cdots & a_{2p}\\
\vdots & & & \vdots\\ a_{n1} & a_{n2} & \cdots & a_{np}\end{pmatrix}.
\]
Le réel $a_{ij}$ (terme de la ligne $i$, colonne $j$) est le \textbf{coefficient} d'indices
$i,j$. On note $\mathcal M_{n,p}(\R)$ l'ensemble des matrices $n\times p$, et $\mathcal M_n(\R)$
celui des matrices carrées d'ordre $n$.
\end{definition}

\begin{definition}[Types usuels]
\begin{itemize}
\item \textbf{Matrice ligne} : un seul ligne, format $1\times p$, ex. $\begin{pmatrix}2 & -1 & 5\end{pmatrix}$.
\item \textbf{Matrice colonne} : une seule colonne, format $n\times 1$, ex. $\begin{pmatrix}3\\0\\-2\end{pmatrix}$.
\item \textbf{Matrice carrée} d'ordre $n$ : autant de lignes que de colonnes ($n=p$). Sa
\textbf{diagonale principale} est formée des $a_{ii}$.
\item \textbf{Matrice nulle} $O$ : tous les coefficients valent $0$.
\item \textbf{Matrice identité} $I_n$ : matrice carrée avec des $1$ sur la diagonale et des $0$ ailleurs,
$I_2=\begin{pmatrix}1&0\\0&1\end{pmatrix}$, $I_3=\begin{pmatrix}1&0&0\\0&1&0\\0&0&1\end{pmatrix}$.
\end{itemize}
\end{definition}

\begin{definition}[Égalité]
Deux matrices $A=(a_{ij})$ et $B=(b_{ij})$ sont \textbf{égales} si elles ont \emph{le même
format} et si $a_{ij}=b_{ij}$ pour tous $i,j$.
\end{definition}

% ═══════════════════════════════════════════════════════════════
\section{Opérations sur les matrices}

\begin{definition}[Somme et produit par un scalaire]
Pour deux matrices $A,B$ de \emph{même format} $n\times p$ et un réel $\lambda$ :
\[
(A+B)_{ij}=a_{ij}+b_{ij},\qquad (\lambda A)_{ij}=\lambda\,a_{ij}.
\]
On ajoute (et on multiplie par un scalaire) \textbf{terme à terme}.
\end{definition}

\begin{propriete}[Règles de calcul (somme)]
Pour $A,B,C$ de même format et $\lambda,\mu\in\R$ :
$A+B=B+A$, \ $(A+B)+C=A+(B+C)$, \ $A+O=A$, \ $A+(-A)=O$, \
$\lambda(A+B)=\lambda A+\lambda B$, \ $(\lambda+\mu)A=\lambda A+\mu A$, \ $\lambda(\mu A)=(\lambda\mu)A$.
\end{propriete}

\begin{exemple}
Avec $A=\begin{pmatrix}1&2\\0&-3\end{pmatrix}$ et $B=\begin{pmatrix}4&-1\\2&5\end{pmatrix}$ :
$A+B=\begin{pmatrix}5&1\\2&2\end{pmatrix}$, \ $2A=\begin{pmatrix}2&4\\0&-6\end{pmatrix}$, \
$2A-B=\begin{pmatrix}-2&5\\-2&-11\end{pmatrix}$.
\end{exemple}

\subsection{Produit matriciel}

\begin{definition}[Produit de deux matrices]
Le produit $A\times B$ n'est défini que si le \textbf{nombre de colonnes de $A$ égale le
nombre de lignes de $B$}. Si $A$ est de format $n\times p$ et $B$ de format $p\times q$, alors
$C=AB$ est de format $n\times q$ et
\[
c_{ij}=\sum_{k=1}^{p}a_{ik}\,b_{kj}=a_{i1}b_{1j}+a_{i2}b_{2j}+\cdots+a_{ip}b_{pj}.
\]
Le coefficient $c_{ij}$ s'obtient en « multipliant la \textbf{ligne $i$ de $A$} par la
\textbf{colonne $j$ de $B$} ».
\end{definition}

\begin{illus}
\begin{tikzpicture}[>=Stealth,scale=1]
% Matrice A (ligne i surlignée)
\node at (0,2.4){\small$A$};
\draw[onbline,line width=0.8pt] (-0.9,1.2) rectangle (0.9,3.0);
\fill[onbo!18] (-0.9,2.2) rectangle (0.9,2.6);
\node at (0,2.4){\small $a_{i1}\;a_{i2}\;a_{i3}$};
\node[onbo,anchor=east] at (-1.0,2.4){\scriptsize ligne $i$};
% Matrice B (colonne j surlignée)
\node at (3.2,3.4){\small$B$};
\draw[onbline,line width=0.8pt] (2.3,1.2) rectangle (4.1,3.0);
\fill[onbblue!16] (2.9,1.2) rectangle (3.3,3.0);
\node[rotate=90] at (3.1,2.1){\small $b_{1j}\;\;b_{2j}\;\;b_{3j}$};
\node[onbblue,anchor=south] at (3.1,3.05){\scriptsize colonne $j$};
% Matrice C résultat
\node at (6.2,2.4){\small$C=AB$};
\draw[onbline,line width=0.8pt] (5.4,1.2) rectangle (7.0,3.0);
\fill[onbgreen!16] (5.9,2.2) rectangle (6.3,2.6);
\filldraw[onbgreen] (6.1,2.4) circle (1.6pt);
\node[onbgreen,anchor=west] at (7.1,2.4){\scriptsize $c_{ij}$};
% flèche
\draw[->,onbmuted,line width=0.7pt] (1.0,2.0) .. controls (3,0.6) and (4.2,0.6) .. (5.9,2.0);
\node[onbmuted,anchor=north] at (3.4,0.55){\scriptsize $c_{ij}=a_{i1}b_{1j}+a_{i2}b_{2j}+a_{i3}b_{3j}$};
\end{tikzpicture}
\fig{Figure — Produit ligne $\times$ colonne : le coefficient $c_{ij}$ combine la ligne $i$ de $A$ et la colonne $j$ de $B$.}
\end{illus}

\begin{exemple}
Soit $A=\begin{pmatrix}1&2\\3&4\end{pmatrix}$ et $B=\begin{pmatrix}5&6\\7&8\end{pmatrix}$.
\[
AB=\begin{pmatrix}1\cdot5+2\cdot7 & 1\cdot6+2\cdot8\\ 3\cdot5+4\cdot7 & 3\cdot6+4\cdot8\end{pmatrix}
=\begin{pmatrix}19&22\\43&50\end{pmatrix}.
\]
On vérifie que $BA=\begin{pmatrix}23&34\\31&46\end{pmatrix}\neq AB$ : le produit n'est
\textbf{pas commutatif}.
\end{exemple}

\begin{propriete}[Règles de calcul (produit)]
Lorsque les formats permettent les produits :
\[
A(BC)=(AB)C,\quad A(B+C)=AB+AC,\quad (A+B)C=AC+BC,\quad \lambda(AB)=(\lambda A)B=A(\lambda B).
\]
Pour toute matrice carrée $A$ d'ordre $n$ : $A\,I_n=I_n\,A=A$ ($I_n$ est l'\textbf{élément neutre}).
\end{propriete}

\begin{aretenir}
Le produit matriciel n'est \textbf{ni commutatif} ($AB\neq BA$ en général), ni intègre :
$AB=O$ \emph{n'entraîne pas} $A=O$ ou $B=O$. Vérifie toujours la \textbf{compatibilité des
formats} ($p$ colonnes de $A$ $=$ $p$ lignes de $B$) avant de poser un produit.
\end{aretenir}

% ═══════════════════════════════════════════════════════════════
\section{Déterminant}

\begin{definition}[Déterminant $2\times2$]
Pour $A=\begin{pmatrix}a&b\\c&d\end{pmatrix}$, le \textbf{déterminant} est
\[
\det A=\begin{vmatrix}a&b\\c&d\end{vmatrix}=ad-bc.
\]
\end{definition}

\begin{definition}[Déterminant $3\times3$ — règle de Sarrus]
Pour $A=\begin{pmatrix}a&b&c\\d&e&f\\g&h&i\end{pmatrix}$ :
\[
\det A=aei+bfg+cdh-ceg-bdi-afh.
\]
On peut aussi développer suivant la première ligne :
$\det A=a\begin{vmatrix}e&f\\h&i\end{vmatrix}-b\begin{vmatrix}d&f\\g&i\end{vmatrix}
+c\begin{vmatrix}d&e\\g&h\end{vmatrix}.$
\end{definition}

\begin{exemple}
$\begin{vmatrix}2&1\\3&4\end{vmatrix}=2\cdot4-1\cdot3=5$. \quad
$\begin{vmatrix}1&2&3\\0&1&4\\5&6&0\end{vmatrix}
=1(1\cdot0-4\cdot6)-2(0\cdot0-4\cdot5)+3(0\cdot6-1\cdot5)=-24+40-15=1.$
\end{exemple}

\begin{theoreme}[Déterminant et produit]
Pour deux matrices carrées de même ordre, $\det(AB)=\det A\times\det B$. En particulier
$A$ est inversible si et seulement si $\det A\neq0$, et alors $\det(A^{-1})=\dfrac1{\det A}$.
\end{theoreme}

% ═══════════════════════════════════════════════════════════════
\section{Matrice inverse}

\begin{definition}[Matrice inversible]
Une matrice carrée $A$ d'ordre $n$ est \textbf{inversible} s'il existe une matrice $A^{-1}$
telle que $A\,A^{-1}=A^{-1}A=I_n$. Cette inverse, si elle existe, est unique.
\end{definition}

\begin{propriete}[Inverse d'une matrice $2\times2$]
$A=\begin{pmatrix}a&b\\c&d\end{pmatrix}$ est inversible si et seulement si $\det A=ad-bc\neq0$,
et alors
\[
A^{-1}=\frac1{ad-bc}\begin{pmatrix}d&-b\\-c&a\end{pmatrix}.
\]
\end{propriete}

\begin{exemple}
$A=\begin{pmatrix}2&1\\3&4\end{pmatrix}$, $\det A=5\neq0$ :
$A^{-1}=\dfrac15\begin{pmatrix}4&-1\\-3&2\end{pmatrix}$. Vérification :
$A\,A^{-1}=\dfrac15\begin{pmatrix}2&1\\3&4\end{pmatrix}\!\begin{pmatrix}4&-1\\-3&2\end{pmatrix}
=\dfrac15\begin{pmatrix}5&0\\0&5\end{pmatrix}=I_2.$
\end{exemple}

\begin{propriete}[Inverse d'une matrice $3\times3$]
Si $\det A\neq0$, l'inverse se calcule par la formule de la \textbf{comatrice} :
\[
A^{-1}=\frac1{\det A}\,{}^{t}\!\big(\operatorname{com}A\big),
\]
où $\operatorname{com}A=\big((-1)^{i+j}\,m_{ij}\big)$ et $m_{ij}$ est le \textbf{mineur}
(déterminant $2\times2$ obtenu en supprimant la ligne $i$ et la colonne $j$).
En pratique on utilise souvent le \textbf{pivot de Gauss} sur $[\,A\mid I_3\,]$.
\end{propriete}

\begin{remarque}
Pour $A,B$ inversibles : $(AB)^{-1}=B^{-1}A^{-1}$ (attention à l'ordre) et $(A^{-1})^{-1}=A$.
\end{remarque}

% ═══════════════════════════════════════════════════════════════
\section{Systèmes linéaires et écriture matricielle}

\begin{definition}[Écriture matricielle $AX=B$]
Le système linéaire
$\left\{\begin{array}{l}a_{11}x+a_{12}y+a_{13}z=b_1\\ a_{21}x+a_{22}y+a_{23}z=b_2\\
a_{31}x+a_{32}y+a_{33}z=b_3\end{array}\right.$
s'écrit $AX=B$ avec
\[
A=\begin{pmatrix}a_{11}&a_{12}&a_{13}\\a_{21}&a_{22}&a_{23}\\a_{31}&a_{32}&a_{33}\end{pmatrix},\quad
X=\begin{pmatrix}x\\y\\z\end{pmatrix},\quad B=\begin{pmatrix}b_1\\b_2\\b_3\end{pmatrix}.
\]
\end{definition}

\begin{theoreme}[Résolution par l'inverse]
Si $A$ est inversible (i.e. $\det A\neq0$), le système $AX=B$ admet une \textbf{unique
solution} donnée par
\[
X=A^{-1}B.
\]
\end{theoreme}

\begin{exemple}
$\left\{\begin{array}{l}2x+y=7\\3x+4y=13\end{array}\right.$ s'écrit $AX=B$ avec
$A=\begin{pmatrix}2&1\\3&4\end{pmatrix}$, $B=\begin{pmatrix}7\\13\end{pmatrix}$. Comme
$A^{-1}=\dfrac15\begin{pmatrix}4&-1\\-3&2\end{pmatrix}$,
$X=A^{-1}B=\dfrac15\begin{pmatrix}28-13\\-21+26\end{pmatrix}=\dfrac15\begin{pmatrix}15\\5\end{pmatrix}
=\begin{pmatrix}3\\1\end{pmatrix}$, soit $x=3,\ y=1$.
\end{exemple}

% ═══════════════════════════════════════════════════════════════
\section{Puissances d'une matrice}

\begin{definition}[Puissance]
Pour une matrice carrée $A$ et $n\in\N^*$, $A^{n}=\underbrace{A\times A\times\cdots\times A}_{n\ \text{facteurs}}$,
et par convention $A^{0}=I_n$.
\end{definition}

\begin{aretenir}
Pour calculer $A^n$, deux grandes techniques au programme :
\begin{enumerate}[label=\arabic*)]
\item \textbf{Conjecture + récurrence} : calculer $A^2$, $A^3$, conjecturer une formule pour $A^n$,
puis la \textbf{démontrer par récurrence}.
\item Écrire $A=I+N$ avec $N$ \textbf{nilpotente} ($N^k=O$) et utiliser le \textbf{binôme de Newton}
(valable car $I$ et $N$ commutent).
\end{enumerate}
\end{aretenir}

\begin{exemple}
Soit $A=\begin{pmatrix}1&1\\0&1\end{pmatrix}$. On calcule
$A^2=\begin{pmatrix}1&2\\0&1\end{pmatrix}$, $A^3=\begin{pmatrix}1&3\\0&1\end{pmatrix}$. On
conjecture $A^n=\begin{pmatrix}1&n\\0&1\end{pmatrix}$, prouvée par récurrence ci-après (méthode~4).
\end{exemple}

% ═══════════════════════════════════════════════════════════════
\section{Matrice d'une transformation du plan}

Dans le plan rapporté à un repère orthonormé, une transformation linéaire associe à un vecteur
$\vec u\,(x,y)$ le vecteur $\vec u\,'(x',y')$ tel que $\begin{pmatrix}x'\\y'\end{pmatrix}
=M\begin{pmatrix}x\\y\end{pmatrix}$, où $M$ est la \textbf{matrice} de la transformation.

\begin{propriete}[Matrices usuelles]
\[
\text{Rotation d'angle }\theta:\ R_\theta=\begin{pmatrix}\cos\theta&-\sin\theta\\\sin\theta&\cos\theta\end{pmatrix};\quad
\text{Homothétie de rapport }k:\ \begin{pmatrix}k&0\\0&k\end{pmatrix};
\]
\[
\text{Sym. axe }(Ox):\begin{pmatrix}1&0\\0&-1\end{pmatrix};\ \
\text{Sym. axe }(Oy):\begin{pmatrix}-1&0\\0&1\end{pmatrix};\ \
\text{Sym. centrale }O:\begin{pmatrix}-1&0\\0&-1\end{pmatrix}.
\]
\end{propriete}

\begin{illus}
\begin{tikzpicture}[>=Stealth,scale=0.95]
\draw[->,onbline] (-0.4,0)--(3.2,0) node[right,onbmuted,font=\scriptsize]{$x$};
\draw[->,onbline] (0,-0.4)--(0,3.2) node[above,onbmuted,font=\scriptsize]{$y$};
% vecteur u = (2.4 , 0.6) approx ; angle ~ 14°
\draw[->,onbo,line width=1.1pt] (0,0)--(2.4,0.7) node[anchor=west,onbo,font=\small]{$\vec u$};
% rotation de 60° de u
\draw[->,onbblue,line width=1.1pt] (0,0)--(0.594,2.444) node[anchor=south,onbblue,font=\small]{$\vec u\,'=R_\theta\vec u$};
% arc d'angle
\draw[onbgreen,line width=0.8pt] (1.1,0.32) arc (16:76:1.15);
\node[onbgreen,font=\small] at (1.35,1.15){$\theta$};
\end{tikzpicture}
\fig{Figure — Une rotation d'angle $\theta$ envoie $\vec u$ sur $R_\theta\vec u$ : $\det R_\theta=\cos^2\theta+\sin^2\theta=1$ (les longueurs sont conservées).}
\end{illus}

\begin{remarque}
La composée de deux transformations correspond au \textbf{produit} de leurs matrices :
$R_\alpha\circ R_\beta$ a pour matrice $R_\alpha R_\beta=R_{\alpha+\beta}$ (on retrouve les
formules d'addition de $\cos$ et $\sin$).
\end{remarque}

% ═══════════════════════════════════════════════════════════════
\section{Méthodes \& savoir-faire}

\methode{Effectuer un produit matriciel}
Vérifier la compatibilité des formats, puis calculer chaque coefficient « ligne $\times$ colonne ».
\begin{exemple}
$\begin{pmatrix}1&-1&2\\0&3&1\end{pmatrix}\begin{pmatrix}2&0\\1&4\\-1&1\end{pmatrix}$ :
format $(2\times3)(3\times2)=2\times2$.
$c_{11}=1\cdot2+(-1)\cdot1+2\cdot(-1)=-1$, $c_{12}=1\cdot0+(-1)\cdot4+2\cdot1=-2$,
$c_{21}=0\cdot2+3\cdot1+1\cdot(-1)=2$, $c_{22}=0\cdot0+3\cdot4+1\cdot1=13$. D'où
$\begin{pmatrix}-1&-2\\2&13\end{pmatrix}$.
\end{exemple}

\methode{Calculer un déterminant et une inverse $2\times2$}
Calculer $\det A=ad-bc$. Si $\det A\neq0$, $A^{-1}=\dfrac1{\det A}\begin{pmatrix}d&-b\\-c&a\end{pmatrix}$.
\begin{exemple}
$A=\begin{pmatrix}3&5\\1&2\end{pmatrix}$ : $\det A=3\cdot2-5\cdot1=1\neq0$, donc
$A^{-1}=\begin{pmatrix}2&-5\\-1&3\end{pmatrix}$. (Vérif. : $A A^{-1}=I_2$.)
\end{exemple}

\methode{Calculer une inverse $3\times3$ par le pivot de Gauss}
Écrire $[\,A\mid I_3\,]$ et appliquer les opérations élémentaires sur les lignes jusqu'à
obtenir $[\,I_3\mid A^{-1}\,]$.
\begin{exemple}
$A=\begin{pmatrix}1&0&2\\0&1&0\\0&0&1\end{pmatrix}$. La seule colonne « à nettoyer » est la 3e :
$L_1\leftarrow L_1-2L_3$ donne $A^{-1}=\begin{pmatrix}1&0&-2\\0&1&0\\0&0&1\end{pmatrix}$.
\end{exemple}

\methode{Calculer $A^n$ par récurrence}
Calculer $A^2,A^3$, conjecturer $A^n$, puis prouver par récurrence en utilisant $A^{n+1}=A^n\cdot A$.
\begin{exemple}
$A=\begin{pmatrix}1&1\\0&1\end{pmatrix}$, conjecture $A^n=\begin{pmatrix}1&n\\0&1\end{pmatrix}$.
\emph{Initialisation} : $A^1=A$, vrai. \emph{Hérédité} : si $A^n=\begin{pmatrix}1&n\\0&1\end{pmatrix}$,
alors $A^{n+1}=A^n A=\begin{pmatrix}1&n\\0&1\end{pmatrix}\begin{pmatrix}1&1\\0&1\end{pmatrix}
=\begin{pmatrix}1&n+1\\0&1\end{pmatrix}$. La formule est donc vraie pour tout $n\ge1$.
\end{exemple}

\methode{Résoudre un système par les matrices}
Écrire $AX=B$, vérifier $\det A\neq0$, calculer $A^{-1}$ et conclure $X=A^{-1}B$.
\begin{exemple}
$\left\{\begin{array}{l}x+2y=5\\3x+5y=11\end{array}\right.$ : $A=\begin{pmatrix}1&2\\3&5\end{pmatrix}$,
$\det A=-1$, $A^{-1}=-\begin{pmatrix}5&-2\\-3&1\end{pmatrix}=\begin{pmatrix}-5&2\\3&-1\end{pmatrix}$.
$X=A^{-1}\begin{pmatrix}5\\11\end{pmatrix}=\begin{pmatrix}-25+22\\15-11\end{pmatrix}=\begin{pmatrix}-3\\4\end{pmatrix}$,
soit $x=-3,\ y=4$.
\end{exemple}

\methode{Déterminer la matrice d'une transformation}
Calculer les images des vecteurs de base $\vec\imath\,(1,0)$ et $\vec\jmath\,(0,1)$ : ce sont
les \emph{colonnes} de la matrice.
\begin{exemple}
Rotation d'angle $\dfrac\pi2$ : $\vec\imath\mapsto(0,1)$, $\vec\jmath\mapsto(-1,0)$, donc
$M=\begin{pmatrix}0&-1\\1&0\end{pmatrix}$, ce qui coïncide avec $R_{\pi/2}$.
\end{exemple}

% ═══════════════════════════════════════════════════════════════
\section{Exercices}

\rubrique{Calcul matriciel}
\exo{}\diff{1} On donne $A=\begin{pmatrix}1&-2\\3&0\end{pmatrix}$ et $B=\begin{pmatrix}4&1\\-1&2\end{pmatrix}$.
Calculer $A+B$, \ $3A-2B$, \ $AB$ et $BA$. Le produit est-il commutatif ?

\exo{}\diff{1} Effectuer le produit
$\begin{pmatrix}2&0&1\\-1&3&2\end{pmatrix}\begin{pmatrix}1&-1\\0&2\\3&1\end{pmatrix}$
(après avoir vérifié sa compatibilité).

\exo{}\diff{2} Soit $N=\begin{pmatrix}0&2\\0&0\end{pmatrix}$. Calculer $N^2$. Que peut-on dire de $N$ ?
En déduire $(I_2+N)^{5}$ à l'aide du binôme.

\rubrique{Déterminant et inverse}
\exo{}\diff{1} Calculer $\det A$ pour $A=\begin{pmatrix}5&3\\2&4\end{pmatrix}$, puis $A^{-1}$.

\exo{}\diff{2} Calculer $\begin{vmatrix}2&-1&0\\1&3&2\\0&1&1\end{vmatrix}$ par la règle de Sarrus.

\exo{}\diff{2} Pour quelles valeurs du réel $m$ la matrice
$A=\begin{pmatrix}m&1\\4&m\end{pmatrix}$ est-elle inversible ? Donner $A^{-1}$ lorsque $m=3$.

\rubrique{Systèmes linéaires}
\exo{}\diff{2} Résoudre par l'écriture matricielle
$\left\{\begin{array}{l}2x+3y=8\\ x-y=-1\end{array}\right.$

\exo{}\diff{3} Soit $A=\begin{pmatrix}1&1&1\\0&1&2\\0&0&1\end{pmatrix}$. Montrer que $A$ est
inversible, calculer $A^{-1}$, puis résoudre
$\left\{\begin{array}{l}x+y+z=6\\ y+2z=8\\ z=3\end{array}\right.$

\rubrique{Puissances et transformations}
\exo{}\diff{2} Soit $A=\begin{pmatrix}3&0\\0&-2\end{pmatrix}$. Conjecturer $A^n$ et le démontrer par récurrence.

\exo{}\diff{3} Déterminer la matrice de la symétrie orthogonale par rapport à la droite $y=x$, en
calculant les images de $\vec\imath$ et $\vec\jmath$. Quel est son déterminant ? Quelle est son inverse ?

\rubrique{Problème type Baccalauréat}
\exo{}\diff{3} On considère la matrice $A=\begin{pmatrix}2&1\\0&3\end{pmatrix}$ et les suites
$(u_n)$, $(v_n)$ définies par $u_0=1$, $v_0=1$ et, pour tout $n\in\N$,
\[
\left\{\begin{array}{l}u_{n+1}=2u_n+v_n\\ v_{n+1}=3v_n\end{array}\right.,\qquad
\text{soit }X_n=\begin{pmatrix}u_n\\v_n\end{pmatrix},\ X_{n+1}=A X_n.
\]
\begin{enumerate}[label=\arabic*)]
\item Calculer $A^2$ et $A^3$.
\item Démontrer par récurrence que pour tout $n\ge1$,
$A^n=\begin{pmatrix}2^n & 3^n-2^n\\ 0 & 3^n\end{pmatrix}$.
\item Justifier que $X_n=A^n X_0$, puis exprimer $u_n$ et $v_n$ en fonction de $n$.
\item Déterminer $\displaystyle\lim_{n\to+\infty}\frac{u_n}{v_n}$.
\end{enumerate}

% ═══════════════════════════════════════════════════════════════
\section{Corrigés}

\corrige{de l'exercice 1}
$A+B=\begin{pmatrix}5&-1\\2&2\end{pmatrix}$. \
$3A-2B=\begin{pmatrix}3&-6\\9&0\end{pmatrix}-\begin{pmatrix}8&2\\-2&4\end{pmatrix}
=\begin{pmatrix}-5&-8\\11&-4\end{pmatrix}$. \
\[
AB=\begin{pmatrix}1\cdot4+(-2)(-1)&1\cdot1+(-2)\cdot2\\ 3\cdot4+0&3\cdot1+0\end{pmatrix}
=\begin{pmatrix}6&-3\\12&3\end{pmatrix},\quad
BA=\begin{pmatrix}4+3&-8+0\\-1+6&2+0\end{pmatrix}=\begin{pmatrix}7&-8\\5&2\end{pmatrix}.
\]
Comme $AB\neq BA$, le produit n'est pas commutatif.

\corrige{de l'exercice 2}
Format $(2\times3)(3\times2)=2\times2$.
$c_{11}=2\cdot1+0\cdot0+1\cdot3=5$, $c_{12}=2(-1)+0\cdot2+1\cdot1=-1$,
$c_{21}=-1\cdot1+3\cdot0+2\cdot3=5$, $c_{22}=-1(-1)+3\cdot2+2\cdot1=9$. Résultat :
$\begin{pmatrix}5&-1\\5&9\end{pmatrix}$.

\corrige{de l'exercice 3}
$N^2=\begin{pmatrix}0&2\\0&0\end{pmatrix}^2=\begin{pmatrix}0&0\\0&0\end{pmatrix}=O$ : $N$ est
\textbf{nilpotente} (d'ordre $2$). Comme $I_2$ et $N$ commutent et $N^2=O$, le binôme donne
$(I_2+N)^5=\sum_{k=0}^{5}\binom5k N^k=I_2+5N=\begin{pmatrix}1&10\\0&1\end{pmatrix}$.

\corrige{de l'exercice 4}
$\det A=5\cdot4-3\cdot2=14\neq0$, donc $A^{-1}=\dfrac1{14}\begin{pmatrix}4&-3\\-2&5\end{pmatrix}$.

\corrige{de l'exercice 5}
Par Sarrus : $aei+bfg+cdh-ceg-bdi-afh$ avec
$(a,b,c)=(2,-1,0)$, $(d,e,f)=(1,3,2)$, $(g,h,i)=(0,1,1)$ :
$2\cdot3\cdot1+(-1)\cdot2\cdot0+0\cdot1\cdot1-0\cdot3\cdot0-(-1)\cdot1\cdot1-2\cdot2\cdot1
=6+0+0-0+1-4=3$.

\corrige{de l'exercice 6}
$\det A=m\cdot m-1\cdot4=m^2-4$. $A$ est inversible $\iff m^2-4\neq0\iff m\neq\pm2$.
Pour $m=3$ : $\det A=5$, $A^{-1}=\dfrac15\begin{pmatrix}3&-1\\-4&3\end{pmatrix}$.

\corrige{de l'exercice 7}
$A=\begin{pmatrix}2&3\\1&-1\end{pmatrix}$, $B=\begin{pmatrix}8\\-1\end{pmatrix}$.
$\det A=2(-1)-3\cdot1=-5\neq0$, $A^{-1}=-\dfrac15\begin{pmatrix}-1&-3\\-1&2\end{pmatrix}
=\dfrac15\begin{pmatrix}1&3\\1&-2\end{pmatrix}$.
$X=A^{-1}B=\dfrac15\begin{pmatrix}8-3\\8+2\end{pmatrix}=\dfrac15\begin{pmatrix}5\\10\end{pmatrix}
=\begin{pmatrix}1\\2\end{pmatrix}$, soit $x=1,\ y=2$.
(Vérification : $2\cdot1+3\cdot2=8$ et $1-2=-1$.)

\corrige{de l'exercice 8}
$A$ est triangulaire à diagonale de $1$ : $\det A=1\neq0$, $A$ est inversible. Par pivot de Gauss
sur $[\,A\mid I_3\,]$ (ou comatrice), on trouve
$A^{-1}=\begin{pmatrix}1&-1&1\\0&1&-2\\0&0&1\end{pmatrix}$.
\emph{Vérification} : $A A^{-1}=\begin{pmatrix}1&-1+1&1-2+1\\0&1&-2+2\\0&0&1\end{pmatrix}=I_3$.\;
Le système s'écrit $AX=B$ avec $B=\begin{pmatrix}6\\8\\3\end{pmatrix}$ :
$X=A^{-1}B=\begin{pmatrix}6-8+3\\8-6\\3\end{pmatrix}=\begin{pmatrix}1\\2\\3\end{pmatrix}$,
soit $x=1,\ y=2,\ z=3$.

\corrige{de l'exercice 9}
$A$ est diagonale : $A^2=\begin{pmatrix}9&0\\0&4\end{pmatrix}$, $A^3=\begin{pmatrix}27&0\\0&-8\end{pmatrix}$.
On conjecture $A^n=\begin{pmatrix}3^n&0\\0&(-2)^n\end{pmatrix}$.
\emph{Initialisation} ($n=1$) : vrai. \emph{Hérédité} : si la formule est vraie au rang $n$,
$A^{n+1}=A^n A=\begin{pmatrix}3^n&0\\0&(-2)^n\end{pmatrix}\begin{pmatrix}3&0\\0&-2\end{pmatrix}
=\begin{pmatrix}3^{n+1}&0\\0&(-2)^{n+1}\end{pmatrix}$. Vrai pour tout $n\ge1$.

\corrige{de l'exercice 10}
La symétrie d'axe $y=x$ échange les coordonnées : $\vec\imath\,(1,0)\mapsto(0,1)$ et
$\vec\jmath\,(0,1)\mapsto(1,0)$. Ces images sont les colonnes :
$M=\begin{pmatrix}0&1\\1&0\end{pmatrix}$. $\det M=0\cdot0-1\cdot1=-1$ (un antidéplacement).
Comme $M^2=\begin{pmatrix}1&0\\0&1\end{pmatrix}=I_2$, $M$ est sa propre inverse : $M^{-1}=M$.

\corrige{du problème}
\textbf{1)} $A^2=\begin{pmatrix}2&1\\0&3\end{pmatrix}^2=\begin{pmatrix}4&5\\0&9\end{pmatrix}$, \
$A^3=A^2A=\begin{pmatrix}4&5\\0&9\end{pmatrix}\begin{pmatrix}2&1\\0&3\end{pmatrix}=\begin{pmatrix}8&19\\0&27\end{pmatrix}$.
On remarque $4=2^2$, $9=3^2$, $5=3^2-2^2$ ; $8=2^3$, $27=3^3$, $19=3^3-2^3$.

\textbf{2)} Conjecture : $A^n=\begin{pmatrix}2^n&3^n-2^n\\0&3^n\end{pmatrix}$.
\emph{Initialisation} ($n=1$) : $\begin{pmatrix}2&3-2\\0&3\end{pmatrix}=\begin{pmatrix}2&1\\0&3\end{pmatrix}=A$, vrai.
\emph{Hérédité} : supposons la formule vraie au rang $n$. Alors
\[
A^{n+1}=A^nA=\begin{pmatrix}2^n&3^n-2^n\\0&3^n\end{pmatrix}\begin{pmatrix}2&1\\0&3\end{pmatrix}
=\begin{pmatrix}2^{n+1}&2^n+3(3^n-2^n)\\0&3^{n+1}\end{pmatrix}.
\]
Or $2^n+3\cdot3^n-3\cdot2^n=3^{n+1}-2^{n+1}$, d'où
$A^{n+1}=\begin{pmatrix}2^{n+1}&3^{n+1}-2^{n+1}\\0&3^{n+1}\end{pmatrix}$. La formule est donc vraie pour tout $n\ge1$.

\textbf{3)} Comme $X_{n+1}=AX_n$, une récurrence immédiate donne $X_n=A^nX_0$ pour tout $n\ge1$
(et c'est aussi vrai pour $n=0$ avec $A^0=I_2$). Avec $X_0=\begin{pmatrix}1\\1\end{pmatrix}$ :
\[
X_n=\begin{pmatrix}2^n&3^n-2^n\\0&3^n\end{pmatrix}\begin{pmatrix}1\\1\end{pmatrix}
=\begin{pmatrix}2^n+3^n-2^n\\3^n\end{pmatrix}=\begin{pmatrix}3^n\\3^n\end{pmatrix}.
\]
Donc $u_n=3^n$ et $v_n=3^n$ pour tout $n\ge0$.
(Contrôle direct : $v_{n+1}=3v_n$ avec $v_0=1$ donne $v_n=3^n$ ; et $u_{n+1}=2u_n+v_n=2u_n+3^n$
avec $u_0=1$ admet bien $u_n=3^n$ pour solution, car $3^{n+1}=2\cdot3^n+3^n$.)

\textbf{4)} $\dfrac{u_n}{v_n}=\dfrac{3^n}{3^n}=1$, donc $\displaystyle\lim_{n\to+\infty}\frac{u_n}{v_n}=1$.
